%-------TITLE PAGE------%
\begin{titlepage}
	\center
	\begin{tikzpicture}[overlay,remember picture]
		\draw [line width=3pt,rounded corners=0pt,]
		($ (current page.north west) + (25mm,-25mm) $)
		rectangle
		($ (current page.south east) + (-15mm,25mm) $);
		\draw [line width=1pt,rounded corners=0pt]
		($ (current page.north west) + (26.5mm,-26.5mm) $)
		rectangle
		($ (current page.south east) + (-16.5mm,26.5mm) $);
	\end{tikzpicture}
	
	{\large \bfseries ĐẠI HỌC QUỐC GIA HÀ NỘI\\ TRƯỜNG ĐẠI HỌC CÔNG NGHỆ}\\[1cm]
	\includegraphics[width=0.25\linewidth]{images/Logo_UET.png}\\[1cm]
	{\Large  \bfseries Trần Hữu Nam}\\[1cm]
% 	{ \LARGE \bfseries  XÂY DỰNG CÔNG CỤ KIỂM THỬ TỰ ĐỘNG }\\[0.05cm]
%     {\LARGE \bfseries CHO CÁC ỨNG DỤNG SỬ DỤNG NGÔN NGỮ}\\[0.2cm]
%      { \LARGE \bfseries   TYPESCRIPT}\\[0.5cm]
% 	\hfill\\[2cm]
	{ \LARGE \bfseries  
        XÂY DỰNG HỆ THỐNG THU THẬP VÀ PHÂN LOẠI TƯ THẾ NGỦ SỬ DỤNG CẢM BIẾN GIA TỐC VÀ MÔ HÌNH HỌC MÁY GỌN NHẸ TRIỂN KHAI TRÊN THIẾT BỊ BIÊN}\\[0.05cm]
	\hfill\\[1cm]
	{\large \bfseries LUẬN VĂN THẠC SĨ}\\	
	% {\large \bfseries Ngành: Kĩ thuật điện tử}	
	\hfill\\[5.3cm]	
	{\large \bfseries HÀ NỘI - 2025}\\	
	\vfill
\end{titlepage}

%-------TITLE PAGE+6hbk,------%
\begin{titlepage}
	\center
	\begin{tikzpicture}[overlay,remember picture]
	\draw [line width=3pt,rounded corners=0pt,]
	($ (current page.north west) + (25mm,-25mm) $)
	rectangle
	($ (current page.south east) + (-15mm,25mm) $);
	\draw [line width=1pt,rounded corners=0pt]
	($ (current page.north west) + (26.5mm,-26.5mm) $)
	rectangle
	($ (current page.south east) + (-16.5mm,26.5mm) $);
	\end{tikzpicture}
	
	{\large \bfseries ĐẠI HỌC QUỐC GIA HÀ NỘI\\ TRƯỜNG ĐẠI HỌC CÔNG NGHỆ}\\[2cm]
% 	\includegraphics[width=0.25\linewidth]{images/Logo_UET.png}\\[1cm]
	
	{\Large  \bfseries Trần Hữu Nam}\\[2cm]		
	{ \LARGE \bfseries  XÂY DỰNG HỆ THỐNG THU THẬP VÀ PHÂN LOẠI TƯ THẾ NGỦ SỬ DỤNG CẢM BIẾN GIA TỐC VÀ MÔ HÌNH HỌC MÁY GỌN NHẸ TRIỂN KHAI TRÊN THIẾT BỊ BIÊN
    }\\[0.05cm]
	\hfill\\[1cm]
	\begin{flushleft}
	{\large \bfseries Ngành: Điện tử viễn thông}\\
	{\large \bfseries Chuyên ngành: Kĩ thuật điện tử}\\
	{\large \bfseries Mã số: 23025029}\\
    
	\end{flushleft}
	\hfill\\[1cm]	
	{\large \bfseries LUẬN VĂN THẠC SĨ}\\
	\hfill\\[1cm]
    
	{\large \bfseries NGƯỜI HƯỚNG DẪN KHOA HỌC: PGS.TS. Mai Anh Tuấn}\\	
    
    
	\begin{flushleft}
% 	{\large \bfseries Cán bộ đồng hướng dẫn: CN. Bùi Quang Cường}\\	
	\end{flushleft}
		\hfill\\[3cm]	
	{\large \bfseries HÀ NỘI - 2025}\\		
	\vfill		
\end{titlepage}

\changefontsizes[16pt]{14pt}
\addtocontents{toc}{\vspace{-1cm}}
\addcontentsline{toc}{chapter}{Lời cam đoan}
\begin{center}
    \textbf{LỜI CAM ĐOAN}
\end{center}
Tôi xin cam đoan: đề tài luận văn thạc sĩ “Xây dựng hệ thống thu thập và phân loại tư thế ngủ sử dụng cảm biến gia tốc và mô hình học máy gọn nhẹ triển khai trên thiết bị biên” là công trình nghiên cứu của tôi dưới sự hướng dẫn của thầy \textbf{PGS. TS. Mai Anh Tuấn} và thầy \textbf{ThS. Trần Ngọc Thái} cùng với sự tham khảo từ những tài liệu đã liệt kê trong mục Tài liệu tham khảo. Tôi không sao chép công trình nghiên cứu của cá nhân khác dưới bất kỳ hình thức nào. Nếu có tôi xin hoàn toàn chịu trách nhiệm.

\begin{flushright}
	\begin{varwidth}{\linewidth}\centering
		Hà Nội, ngày \space\space\space\space tháng  \space\space\space\space năm 2025\\
		Học viên\\[2cm]
		Trần Hữu Nam
	\end{varwidth}
\end{flushright}

\newpage

\addcontentsline{toc}{chapter}{Lời cảm ơn}
\begin{center}
    \textbf{LỜI CẢM ƠN}
\end{center}

Lời đầu tiên, tôi xin gửi lời cảm ơn đến thầy PGS.TS. Mai Anh Tuấn và thầy ThS. Trần Ngọc Thái vì đã tận tình hướng dẫn, truyền đạt kiến thức cho tôi trong suốt quá trình học tập và thực hiện đề tài. 
Tôi xin cảm ơn tập thể thầy, cô khoa “Điện tử Viễn thông”,  Trường Đại học Công nghệ - ĐHQGHN, đã giảng dạy tôi trong quá trình tôi học tập tại trường. Tôi cũng xin cảm ơn các anh chị ở Bộ môn Công nghệ Vi cơ Điện tử và Kỹ thuật Máy tính đã tạo điều kiện giúp đỡ, chỉ bảo tôi trong thời gian làm khóa luận. 
Cuối cùng, tôi xin cảm ơn bố mẹ, gia đình cũng như bạn bè, tập thể lớp K30 đã luôn đồng hành, chia sẻ và động viên tôi suốt thời gian qua.


\newpage
\addcontentsline{toc}{chapter}{Tóm tắt}
\begin{center}
    \textbf{TÓM TẮT}
\end{center}
% \changefontsizes[16pt]{12pt}
\textit{\textbf{Tóm tắt: }} 
Ngưng thở tắc nghẽn khi ngủ (Obstructive Sleep Apnea - OSA) là tình trạng rối loạn hô hấp khi ngủ thường gặp, được đặc trưng bởi những cơn ngưng/giảm thở ngắn, lặp lại trong khi ngủ, gây gián đoạn giấc ngủ bởi hành vi thức giấc thường xuyên và giảm oxy máu ngắt quãng\cite{Epstein2009}. Nghiên cứu\cite{Salari2025} đã phân tích dữ liệu từ 15 nghiên cứu với tổng cộng 42.924 người. Kết quả cho thấy tỷ lệ mắc chứng buồn ngủ ban ngày quá mức (EDS) ở bệnh nhân ngưng thở tắc nghẽn khi ngủ (OSA) trên toàn cầu là 39,9\% (khoảng tin cậy 95\%: 34,4\%–45,7\%).

Tư thế ngủ đóng vai trò quan trọng trong việc khởi phát và làm trầm trọng thêm các 
triệu chứng của hội chứng ngưng thở khi ngủ (Obstructive Sleep Apnea – OSA), 
đặc biệt là dạng OSA theo tư thế (positional OSA). 
Các nghiên cứu cho thấy tần suất ngưng thở tăng rõ rệt khi bệnh nhân nằm ngửa, 
so với các tư thế khác như nằm nghiêng trái hoặc phải. 
Trong bối cảnh đó, việc theo dõi chính xác tư thế ngủ trong thời gian thực có thể cung cấp thông tin quan trọng phục vụ chẩn đoán sớm, 
đánh giá nguy cơ và cải thiện tình trạng OSA bằng việc thay đổi tư thế ngủ.

Trong bối cảnh công nghệ chế tạo ngày càng phát triển mạnh mẽ, việc thu nhỏ hóa và tối ưu hiệu suất hoạt động của các loại cảm biến, pin đã trở thành một yếu tố then chốt, mang tính quyết định trong quá trình tích hợp chúng vào các thiết bị điện tử có kích thước nhỏ gọn. Không chỉ góp phần nâng cao độ chính xác trong việc đo lường các thông số sinh lý quan trọng, nó còn giúp giảm thiểu kích thước thiết bị, tăng tính di động và mở rộng khả năng ứng dụng trong nhiều lĩnh vực, đặc biệt nổi bật trong y học cá thể hóa.

Bên cạnh đó, sự phát triển vượt bậc của trí tuệ nhân tạo (AI) đóng vai trò như một đòn bẩy quan trọng trong việc nâng cao hiệu quả khai thác dữ liệu cảm biến. AI không chỉ tối ưu hóa quy trình xử lý và phân tích dữ liệu mà còn thúc đẩy khả năng nhận diện mẫu, phân cụm, dự đoán và đưa ra quyết định mang tính hỗ trợ lâm sàng. Trong lĩnh vực chẩn đoán hội chứng ngưng thở khi ngủ, sự phối hợp giữa cảm biến và thuật toán học máy không chỉ đảm bảo độ chính xác cao trong thu thập dữ liệu mà còn mở ra khả năng phân tích chuyên sâu về các yếu tố sinh lý, phục vụ quá trình đánh giá y tế toàn diện.
Đặc biệt, sự nổi lên của lĩnh vực học máy triển khai trực tiếp trên các vi điều khiển hoặc thiết bị biên có tài nguyên hạn chế (Tiny Machine Learning - TinyML) đã đánh dấu bước tiến quan trọng trong việc hiện thực hóa các hệ thống giám sát sức khỏe thuận tiện, tiết kiệm năng lượng, chi phí thấp và có khả năng hoạt động độc lập không phụ thuộc vào kết nối mạng hoặc nền tảng đám mây. TinyML không chỉ giúp thu hẹp khoảng cách giữa nghiên cứu học thuật và ứng dụng thực tiễn, mà còn mở ra cơ hội phổ cập các công nghệ theo dõi sức khỏe đến cộng đồng.


Trên cơ sở đó, nghiên cứu này tập trung phát triển một hệ thống theo dõi tình trạng bệnh nhân trong khi ngủ, 
sử dụng cảm biến gia tốc kết hợp với thuật toán học máy nhằm phân tích tư thế ngủ. 
Hệ thống được thiết kế bao gồm hai giai đoạn chính: 
đầu tiên là xây dựng ứng dụng di động phục vụ thu thập và tiền xử lý dữ liệu cảm biến theo thời gian thực; 
tiếp theo là triển khai mô hình học máy gọn nhẹ lên vi điều khiển, 
hướng tới khả năng giám sát liên tục ngay trên thiết bị biên. 
Giải pháp này cho phép thực hiện đánh giá tư thế ngủ với độ chính xác cao, độ trễ thấp,
đồng thời đảm bảo tính cơ động, bảo mật và khả năng hoạt động độc lập không phụ thuộc vào hạ tầng mạng. 
Việc tích hợp toàn bộ quy trình từ thu thập đến phân tích dữ liệu ngay tại thiết bị hứa hẹn nâng cao hiệu quả phát hiện sớm và 
hỗ trợ điều trị các bệnh lý hô hấp liên quan đến giấc ngủ một cách kịp thời và bền vững.



\vspace{-0.5cm}
\begin{flushleft}
  \textit{\textbf{Từ khóa: } cảm biến gia tốc, học máy, ngưng thở tắc nghẽn khi ngủ,  Tiny ML}
\end{flushleft}
% \changefontsizes[16pt]{13pt}



\changefontsizes[16pt]{13pt}
